\documentclass[11pt,a4paper]{article}
\usepackage{planos-ufmt}
\begin{document}
\CapaPlano{Teorema de Existência e Unicidade para EDOs}{UFMT Várzea Grande, área Matemática Aplicada/Probabilidade, prova didática 2025}{EDOs; graduação com espaços métricos, norma supremo e integração.}
\section{Objetivos e estratégia}
Demonstrar uma versão local do teorema de Picard--Lindelöf usando contração, explicar por que a garantia é local e resolver um exemplo por iterações. Método: transformar problema diferencial em equação integral, estimar operador e usar ponto fixo.
\Cronograma{Apresentar PVI $x'=f(t,x)$ e perguntar se continuidade basta para unicidade.}{Enunciar hipótese local de continuidade e Lipschitz em $x$; construir retângulo e espaço de funções.}{Demonstrar que o operador integral é uma contração numa bola fechada; concluir existência e unicidade local.}{Calcular iterações de Picard para $x'=x$, $x(0)=1$; contrastar com $x'=\sqrt{|x|}$, $x(0)=0$.}{Concluir diferença entre existência local e prolongamento global.}
\section{Demonstração por ponto fixo}
\begin{teorema}[Picard--Lindelöf, versão local]
Se $f$ é contínua no retângulo $R=[t_0-a,t_0+a]\times[x_0-b,x_0+b]$ e Lipschitz em $x$ ali, com constante $L$, então existe $T>0$ tal que o PVI $x'=f(t,x)$, $x(t_0)=x_0$, tem solução única em $[t_0-T,t_0+T]$ que permanece no retângulo.
\end{teorema}
\begin{proof}
Seja $M=\max_R|f|$ e escolha $T>0$ com $T\le a$, $TM\le b$ e $TL<1$; quando $M=0$ ou $L=0$, as respectivas condições são automáticas. Considere o conjunto fechado
\[
\mathcal B=\{u\in C([t_0-T,t_0+T]):\|u-x_0\|_\infty\le b\},
\]
completo para a norma supremo. Defina
\[
(\Phi u)(t)=x_0+\int_{t_0}^t f(s,u(s))\,ds.
\]
Para $u\in\mathcal B$, vale $|\Phi u(t)-x_0|\le M|t-t_0|\le TM\le b$; então $\Phi$ leva $\mathcal B$ nela mesma. Ademais,
\[
\|\Phi u-\Phi v\|_\infty
\le LT\|u-v\|_\infty,
\]
inclusive para $t<t_0$, usando a integral orientada. Como $LT<1$, Banach fornece único ponto fixo $x\in\mathcal B$. Pelo Teorema Fundamental do Cálculo, a equação integral implica $x'=f(t,x)$ e $x(t_0)=x_0$. Reciprocamente qualquer solução que permaneça no retângulo e no intervalo escolhido é ponto fixo, logo coincide com $x$.
\end{proof}
\section{Exemplos resolvidos e limites}
Para $x'=x$, $x(0)=1$, inicie $u_0(t)=1$ e defina $u_{n+1}(t)=1+\int_0^t u_n(s)\,ds$. Obtém-se
\[
u_1=1+t,\quad u_2=1+t+\tfrac{t^2}{2!},\quad
u_n=\sum_{k=0}^{n}\tfrac{t^k}{k!}\longrightarrow e^t.
\]
\textbf{Por que Lipschitz importa?} No PVI $x'=\sqrt{|x|}$, $x(0)=0$, tanto $x(t)\equiv0$ quanto $x(t)=t^2/4$ para $t\ge0$, prolongado por $0$ para $t\le0$, são soluções: unicidade pode falhar com $f$ contínua que não é localmente Lipschitz em $x=0$.
\Fechamento{Esquerda: retângulo e $\mathcal B$; centro: estimativa $LT$; direita: série de Picard e contraexemplo.}{Local significa várias soluções? Não: significa garantia numa vizinhança, podendo haver prolongamento único posterior sob hipóteses.}
\section{Autoavaliação específica}
\begin{itemize}
\item Provei que $\Phi(\mathcal B)\subseteq\mathcal B$ e que $\Phi$ é contração?
\item Expliquei o papel de cada limitação sobre $T$?
\item Diferenciei unicidade local de existência global?
\end{itemize}
\RubricaUFMT
\section{Referências}
J. Sotomayor, \emph{Equações Diferenciais Ordinárias}; D. G. Figueiredo e A. F. Neves, \emph{Equações Diferenciais Aplicadas}.

\end{document}
